% --- 값과 거리: 점과 직선 사이 거리 공식 → 분모 = 1 이면 거리 = |값|
\begin{frame}{왜 $w_1^2 + w_2^2 = 1$ 인가: 점과 직선 사이 거리 공식}
\vfill
\begin{columns}[c]
\begin{column}{0.42\textwidth}\centering\includegraphics[width=\linewidth]{sv_dist.png}\end{column}
\begin{column}{0.56\textwidth}
{\large 점과 직선 사이 거리 공식을}\\[0.2em]
{\large 직선 $w_1x_1 + w_2x_2 + b = 0$ 에 쓰면}\\[0.4em]
{\Large $\displaystyle d = \frac{|w_1 x_1 + w_2 x_2 + b|}{\textcolor{ksared}{\sqrt{w_1^2 + w_2^2}}}$}\\[0.6em]
{\large 분모를 1 로 맞추면 \;$\Rightarrow$\; \kb{$d$ = |식의 값|}}\\[0.3em]
{\normalsize 그림: 값 1.41 $\div \sqrt2$ = 거리 1}
\end{column}
\end{columns}
\vfill
\note{점과 직선 사이의 거리 공식을 떠올려 봅시다. 점과 직선 a x1 더하기 b x2 더하기 c가 0인 직선 사이 거리는, 점을 식에 넣은 값의 절댓값을 루트 a제곱 더하기 b제곱으로 나눈 것입니다. 초평면 식으로 쓰면 분자는 w1 x1 더하기 w2 x2 더하기 b의 절댓값, 분모는 루트 w1제곱 더하기 w2제곱입니다. 그러니 분모가 1이 되도록, 즉 w1제곱 더하기 w2제곱을 1로 맞춰 두면, 식에 점을 넣은 값이 그대로 거리가 됩니다. 앞 쪽 식의 둘째 줄 조건이 바로 이것입니다. 그림에서 x1 더하기 x2 빼기 3에 거리 1인 점을 넣으면 값이 1.41인데, 루트 2로 나누면 거리 1입니다. 처음부터 식을 루트 2로 나눠 0.71, 0.71, 마이너스 2.12로 쓰면 값이 바로 1.00, 거리와 같습니다. 이 공식이 왜 성립하는지는 다음 참고 쪽에 있습니다.}
\end{frame}

